\chapter{Games, Auctions and Information}\label{ch:qm:games-auctions-and-information}

In November 1998 the United States Treasury moved every auction of marketable securities to a single clearing price: all winners pay the highest accepted yield, instead of each paying its own bid.
It had run an experiment with two- and five-year notes since September 1992, and its own study of the experiment found that dealers bid more aggressively under the single price, that awards were less
concentrated, and that the auction's yield spread over the when-issued market roughly halved for two-year notes (from 0.41 to 0.22 of a basis point), a difference that was not statistically
significant because results were more volatile from auction to auction. A dealer bids differently in the two formats. The theory of games and auctions says how, and whether the Treasury should have
expected to raise more. This chapter develops that theory, and ends with the information theory of a bettor who buys a tip.

\section{Games and equilibrium}

\begin{definition}[Normal-form game, Nash equilibrium, mixed strategy, zero-sum game]\label{def:qm:games-auctions-and-information:game}
A \emph{normal-form game}\index{normal-form game} is a set of players, a set of strategies for each and a payoff to each player for every profile of strategies. A \emph{mixed strategy}\index{mixed strategy} is a
probability distribution over a player's strategies. A \emph{Nash equilibrium}\index{Nash equilibrium} is a profile of (mixed) strategies in which no player gains by deviating alone. A \emph{zero-sum
game}\index{zero-sum game} is a two-player game whose payoffs sum to zero; with a payoff matrix $A$ for the row player, mixed strategies $x$ and $y$ give her $x^\top Ay$.
\end{definition}

\begin{theorem}[Existence of equilibria; minimax]\label{thm:qm:games-auctions-and-information:minimax}
Every finite game has a \omterm{def:qm:games-auctions-and-information:game}{Nash equilibrium} in mixed strategies (Nash, 1950; stated). In a finite \omterm{def:qm:games-auctions-and-information:game}{zero-sum game}, $\max_x\min_yx^\top Ay = \min_y\max_xx^\top Ay$ (von Neumann, 1928); the common value is the value
of the game, and the optimal strategies form the equilibria.
\end{theorem}

\begin{proof}[Idea]
Existence follows from Kakutani's fixed-point theorem applied to the best-response correspondence. For the minimax theorem, the row player's problem $\max\{v : A^\top x \ge v\mathbf 1, \mathbf 1^\top x = 1, x \ge
0\}$ is a \omterm{def:qm:convex-optimisation:classes}{linear programme} whose dual is the column player's problem, and \omterm{def:qm:convex-optimisation:strong}{strong duality} (chapter~23) equates their values.
\end{proof}

In weighted rock--paper--scissors (rock beats scissors for 1, scissors beat paper for 1, paper beats rock for 2), the equilibrium is not uniform: each player plays rock, paper and scissors with probabilities
$(\frac14, \frac12, \frac14)$, and the value is zero. Two players who each run multiplicative weights for 20\,000 rounds, reinforcing the actions that would have paid, average to $(0.250, 0.502, 0.248)$ with a
duality gap of 0.005: the learning dynamics find the equilibrium neither player computed. Nash equilibria of non-zero-sum games are harder to compute and need not be unique; a trading desk meets them in
quoting games (Book 2) and in the auctions below.

\section{Bayesian games}

\begin{definition}[Bayesian game, Bayes--Nash equilibrium]\label{def:qm:games-auctions-and-information:bayes}
A \emph{Bayesian game}\index{Bayesian game} (Harsanyi, 1967--68) adds to a game a type for each player, drawn from a common prior and known only to its owner; a strategy maps types to actions. A
\emph{Bayes--Nash equilibrium}\index{Bayes--Nash equilibrium} is a strategy profile in which each type's action maximises its expected payoff given its beliefs about the others' types and their strategies.
\end{definition}

\begin{definition}[Private and common values, first- and second-price auctions]\label{def:qm:games-auctions-and-information:auctions}
In a \emph{private-value auction}\index{private-value auction} each bidder knows its own value for the object; in a \emph{common-value auction}\index{common-value auction} the value is the same for all and each
bidder sees only a signal of it. In a \emph{first-price auction}\index{first-price auction} sealed bids are opened and the highest bidder pays its bid; in a \emph{second-price auction}\index{second-price auction}
(Vickrey, 1961) the highest bidder pays the second-highest bid. \emph{Bid shading}\index{bid shading} is bidding below one's value.
\end{definition}

\begin{proposition}[Equilibrium bids]\label{prop:qm:games-auctions-and-information:bids}
With $n$ bidders whose values are independent with distribution function $F$ on $[0, 1]$: in the \omterm{def:qm:games-auctions-and-information:auctions}{second-price auction}, bidding one's value is weakly dominant; in the \omterm{def:qm:games-auctions-and-information:auctions}{first-price auction}, the symmetric \omterm{def:qm:games-auctions-and-information:bayes}{Bayes--Nash
equilibrium} is $b(v) = v - \int_0^vF(x)^{n-1}dx/F(v)^{n-1}$, which for uniform values is $b(v) = \frac{n-1}nv$.
\end{proposition}

\begin{proof}
Second price: bidding $b \ne v$ changes the outcome only when the highest rival bid $h$ lies between $b$ and $v$; if $b > h > v$ the bidder wins and pays $h > v$, a loss, and if $v > h > b$ it forgoes the
profit $v - h$. First price: if rivals use an increasing $b$, a bidder with value $v$ who bids $b(w)$ wins with probability $F(w)^{n-1}$ and earns $(v - b(w))F(w)^{n-1}$. The first-order condition at $w = v$ is
$\frac{d}{dv}(b(v)F(v)^{n-1}) = v\frac{d}{dv}F(v)^{n-1}$, and integrating by parts from $b(0) = 0$ gives the formula.
\end{proof}

With five bidders and uniform values each shades its bid by exactly one fifth. The equilibrium is a best response bidder by bidder: a bidder with value 0.8 facing four rivals who bid $\frac45$ of their values
maximises its expected profit at a bid of 0.64 (\cref{fig:qm:games-auctions-and-information:br}). The shading is the price of not paying the second-highest bid: it recovers exactly the expected gap between
one's own value and the next one below.

\begin{omfigure}
\begin{tikzpicture}
\begin{axis}[width=10cm,height=4.8cm,xlabel={bid},ylabel={expected profit},
  tick label style={font=\small},label style={font=\small},xmin=0,xmax=0.8,ymin=0,ymax=0.08,grid=major,grid style={gray!25},
  scaled y ticks=false,yticklabel style={/pgf/number format/fixed,/pgf/number format/precision=2},table/col sep=comma]
\addplot[omDef,very thick] table[x=bid,y=payoff]{figdata/methods/29-games-auctions-and-information/bestresponse.csv};
\draw[omThm,dashed,thick] (axis cs:0.64,0) -- (axis cs:0.64,0.08) node[pos=0.9,left,font=\footnotesize]{$b(0.8) = 0.64$};
\end{axis}
\end{tikzpicture}

\omcaption{Expected profit $(v - b)\,\P(\text{win})$ of a bidder with value 0.8 in a \omterm{def:qm:games-auctions-and-information:auctions}{first-price auction} against four rivals with uniform values who bid $\frac45$ of their values: the best response is the
equilibrium bid. Data: the chapter's tutorial.}\label{fig:qm:games-auctions-and-information:br}
\end{omfigure}

\section{Auction formats and revenue}

\begin{definition}[English and Dutch auctions, reserve price]\label{def:qm:games-auctions-and-information:open}
An \emph{English auction}\index{English auction} raises the price until one bidder remains, who pays the price at which the last rival dropped out. A \emph{Dutch auction}\index{Dutch auction} lowers the price
until a bidder accepts it. A \emph{reserve price}\index{reserve price} is the lowest price at which the seller will sell.
\end{definition}

The \omterm{def:qm:games-auctions-and-information:open}{English auction} is strategically a \omterm{def:qm:games-auctions-and-information:auctions}{second-price auction} with private values (stay in until the price reaches your value), and the \omterm{def:qm:games-auctions-and-information:open}{Dutch auction} is a \omterm{def:qm:games-auctions-and-information:auctions}{first-price auction} (choose the price at which to stop).

\begin{theorem}[Revenue equivalence]\label{thm:qm:games-auctions-and-information:re}
With independent private values drawn from a common continuous distribution and risk-neutral bidders, every auction in which the highest value wins in equilibrium and a bidder with the lowest possible value
expects zero surplus yields the same expected revenue (Vickrey, 1961; Myerson, 1981; Riley and Samuelson, 1981).
\end{theorem}

\begin{proof}[Idea]
In any such equilibrium, a bidder's expected surplus $U(v)$ satisfies the envelope condition $U'(v) = \P(\text{win} \mid v) = F(v)^{n-1}$, so $U(v) = U(0) + \int_0^vF^{n-1}$ is fixed by the allocation rule. Its
expected payment, $vF(v)^{n-1} - U(v)$, is then the same in every format, and so is the seller's revenue, the sum of the expected payments.
\end{proof}

With five bidders and uniform values, the expected revenue of every format is $\frac{n-1}{n+1} = \frac23$. A million simulated auctions of each give 0.6665 (first price), 0.6665 (second price), 0.6715 for
an English clock that rises by ticks of 0.01 and 0.6614 for a Dutch clock that falls by ticks of 0.01, the last two off by half a tick in opposite directions. The first-price winner pays its shaded bid, the
second-price winner the second value; the averages agree, the distributions do not.

\begin{proposition}[Optimal reserve price]\label{prop:qm:games-auctions-and-information:reserve}
For regular value distributions, the revenue-maximising auction sets a reserve $r^*$ solving $r - (1 - F(r))/f(r) = v_0$, the seller's own value, independent of the number of bidders (Myerson, 1981;
stated). For uniform values and $v_0 = 0$, $r^* = \frac12$, and with $n$ bidders the expected revenue is $\frac{n-1}{n+1} - \frac{2n}{n+1}r^{n+1} + r^n$.
\end{proposition}

A reserve raises revenue by sometimes not selling. With two bidders the optimal reserve of 0.5 raises the expected revenue from $\frac13$ to $\frac5{12}$, by 25\%, and leaves the object unsold a quarter of the time;
with five bidders it raises it from 0.6667 to 0.6719, by 0.8\%, and leaves it unsold 3.1\% of the time (\cref{fig:qm:games-auctions-and-information:reserve}): competition does the reserve's work.

\begin{omfigure}
\begin{tikzpicture}
\begin{axis}[width=10cm,height=5cm,xlabel={reserve price $r$},ylabel={expected revenue},
  tick label style={font=\small},label style={font=\small},xmin=0,xmax=1,ymin=0,ymax=0.75,grid=major,grid style={gray!25},
  legend cell align=left,legend style={font=\footnotesize,at={(0.03,0.03)},anchor=south west,fill=white,draw=none},table/col sep=comma]
\addplot[omThm,very thick] table[x=r,y=n2]{figdata/methods/29-games-auctions-and-information/reserve.csv};
\addplot[omDef,very thick] table[x=r,y=n5]{figdata/methods/29-games-auctions-and-information/reserve.csv};
\draw[black,dashed] (axis cs:0.5,0) -- (axis cs:0.5,0.75);
\legend{two bidders,five bidders}
\end{axis}
\end{tikzpicture}

\omcaption{Expected revenue of a standard auction with uniform private values against the \omterm{def:qm:games-auctions-and-information:open}{reserve price}; both curves peak at $r^* = \frac12$. Data: the chapter's tutorial (closed form, checked by
simulation).}\label{fig:qm:games-auctions-and-information:reserve}
\end{omfigure}

\begin{definition}[Pay-as-bid and discriminatory auctions]\label{def:qm:games-auctions-and-information:pab}
A \emph{pay-as-bid auction}\index{pay-as-bid auction}, or \emph{discriminatory auction}\index{discriminatory auction}, sells several units to the highest bids, each winner paying its own bid; it is the multi-unit
\omterm{def:qm:games-auctions-and-information:auctions}{first-price auction}, as the uniform-price auction of Book 2 is the multi-unit \omterm{def:qm:games-auctions-and-information:auctions}{second-price auction}.
\end{definition}

The Treasury's switch can now be modelled. Take five dealers, each wanting one of two units, with uniform private values. In the uniform-price auction each bids its value, and revenue is twice the third-highest
value, $2 \times \frac36 = 1$ in expectation. In the \omterm{def:qm:games-auctions-and-information:pab}{pay-as-bid auction} a dealer with value $v$ bids the expected value of the second-highest of the other four given that it is below $v$, $b(v) = v(3 - 2.4v)/(4 -
3v)$: 0.36 at a value of 0.5, 0.54 at 0.8, 0.60 at 1. By revenue equivalence expected revenue is again 1; 400\,000 simulated auctions give 0.9994 and 0.9999. What differs is everything else: under the uniform price the
bids are spread twice as widely (the gap between the two winning bids averages 0.167 against 0.087) and the revenue varies more than twice as much from auction to auction (standard deviation 0.378 against 0.160;
\cref{fig:qm:games-auctions-and-information:multi}), the two features the Treasury's study reported. With unit demand the model cannot favour either format on average; real dealers demand many units and can
shade their bids in a uniform-price auction too (demand reduction), which is why the question was empirical.

\begin{omfigure}
\begin{tikzpicture}
\begin{axis}[width=10cm,height=5cm,xlabel={revenue of one auction},ylabel={density},
  tick label style={font=\small},label style={font=\small},xmin=0,xmax=2,ymin=0,ymax=3.8,grid=major,grid style={gray!25},
  legend cell align=left,legend style={font=\footnotesize,at={(0.97,0.97)},anchor=north east,fill=white,draw=none},table/col sep=comma]
\addplot[omThm,very thick] table[x=revenue,y=uniform]{figdata/methods/29-games-auctions-and-information/multiunit.csv};
\addplot[omDef,very thick] table[x=revenue,y=payasbid]{figdata/methods/29-games-auctions-and-information/multiunit.csv};
\legend{uniform price,pay-as-bid}
\end{axis}
\end{tikzpicture}

\omcaption{Distribution of revenue in 400\,000 simulated auctions of two units to five unit-demand bidders with uniform private values, under the uniform-price and pay-as-bid formats in equilibrium: the same
mean, 1, and very different spreads. Data: the chapter's tutorial.}\label{fig:qm:games-auctions-and-information:multi}
\end{omfigure}

Common values change the logic. If the value $V$ is common and each of five bidders sees $V$ plus uniform noise on $\pm0.1$, the highest signal overstates $V$ by $0.1 \times \frac46 = 0.067$ on average, and a bidder
who bids its signal in a \omterm{def:qm:games-auctions-and-information:auctions}{first-price auction} loses that much when it wins, in 97\% of its wins: the winner's curse of Book 2. Shading by the expected overstatement brings the winners' profit to zero. With
affiliated values the formats no longer raise the same revenue: the \omterm{def:qm:games-auctions-and-information:open}{English auction}, which reveals the dropouts' information, raises at least as much as the \omterm{def:qm:games-auctions-and-information:auctions}{second-price auction}, which raises at least as much as the
\omterm{def:qm:games-auctions-and-information:auctions}{first-price auction} (Milgrom and Weber, 1982).

\section{Entropy and optimal betting}

\begin{definition}[Entropy, mutual information]\label{def:qm:games-auctions-and-information:entropy}
The \emph{entropy}\index{entropy} of a discrete random variable $X$ is $H(X) = -\sum_xp(x)\log_2p(x)$ bits (Shannon, 1948). The \emph{mutual information}\index{mutual information} of $X$ and $Y$ is $I(X; Y) =
\sum_{x,y}p(x, y)\log_2\frac{p(x, y)}{p(x)p(y)} = H(X) - H(X \mid Y)$, the \omterm{def:qm:estimation:kl}{Kullback--Leibler divergence} of the joint law from the product of its marginals.
\end{definition}

\begin{theorem}[Kelly: the value of side information]\label{thm:qm:games-auctions-and-information:kelly}
In a horse race with win probabilities $p(x)$ and decimal odds $o(x)$, a bettor who stakes its whole wealth in proportions $b(x)$ has doubling rate $W(b) = \sum_xp(x)\log_2(b(x)o(x))$, maximised by $b = p$ (the
Kelly criterion of Book 2). With a signal $Y$ observed before each race, the optimal doubling rate rises by exactly $I(X; Y)$ bits per race, whatever the odds (Kelly, 1956).
\end{theorem}

\begin{proof}
$W(b) - W(p) = \sum_xp(x)\log_2(b(x)/p(x)) = -D(p\,\|\,b) \le 0$, so $b = p$ is optimal and $W^* = \sum_xp(x)\log_2(p(x)o(x))$. With the signal, the optimal stakes are $p(x \mid y)$ and
\[W^*_Y = \sum_{x,y}p(x, y)\log_2(p(x \mid y)o(x)).\] The odds cancel in the difference: $W^*_Y - W^* = \sum_{x,y}p(x, y)\log_2\frac{p(x \mid y)}{p(x)} = I(X; Y)$.
\end{proof}

A four-horse race with probabilities $(0.4, 0.3, 0.2, 0.1)$ and decimal odds $(2.2, 3, 4.5, 9)$, where the track keeps about 11\% (the implied probabilities sum to 1.12), has \omterm{def:qm:games-auctions-and-information:entropy}{entropy} 1.846 bits. The best bettor
loses 0.165 bits a race: its wealth halves every six races or so. A tip that names the winner with probability 0.6 (and one of the others at random otherwise) carries 0.363 bits, and turns the doubling rate to $+0.198$;
200\,000 simulated races give a gain of $0.363 \pm 0.002$. A tip right with probability $\frac14$ carries nothing, and a perfect tip carries the whole \omterm{def:qm:games-auctions-and-information:entropy}{entropy} (\cref{fig:qm:games-auctions-and-information:kelly}).
Information has a price in bits, and the bits are exactly the growth.

\begin{omfigure}
\begin{tikzpicture}
\begin{axis}[width=10cm,height=5cm,xlabel={probability that the tip names the winner},ylabel={bits per race},
  tick label style={font=\small},label style={font=\small},xmin=0.2,xmax=1.02,ymin=0,ymax=2,grid=major,grid style={gray!25},
  legend cell align=left,legend style={font=\footnotesize,at={(0.03,0.97)},anchor=north west,fill=white,draw=none},table/col sep=comma]
\addplot[omDef,very thick] table[x=q,y=mi]{figdata/methods/29-games-auctions-and-information/kelly.csv};
\addplot[only marks,mark=*,mark size=1.6pt,omThm] table[x=q,y=sim]{figdata/methods/29-games-auctions-and-information/kelly.csv};
\legend{mutual information $I(X; Y)$,simulated gain in doubling rate}
\end{axis}
\end{tikzpicture}

\omcaption{Gain in the Kelly bettor's doubling rate from a noisy tip, against the tip's accuracy, in a four-horse race: the \omterm{def:qm:games-auctions-and-information:entropy}{mutual information} and the gain measured over 200\,000 simulated races for each accuracy.
Data: the chapter's tutorial, seeded.}\label{fig:qm:games-auctions-and-information:kelly}
\end{omfigure}

\section{Tutorial: the Treasury's switch}\label{tut:qm:games-auctions-and-information:switch}

\begin{tutorial}
\textbf{Goal.} Check revenue equivalence in four single-unit formats and two multi-unit formats, price a reserve, show the winner's curse, and measure the value of a tip in bits. \textbf{End state:} the
figures and the named numbers of the weekend problem.
\begin{tutsteps}
\item \textbf{Equilibrium bids.} First-price (closed form and general), and pay-as-bid with unit demand.
\omcode{code/firm/bidding/firm_bidding.py}{26}{58}{Equilibrium bid functions (Python).}
\item \textbf{The value of a tip.} The joint law of winner and tip, the \omterm{def:qm:games-auctions-and-information:entropy}{mutual information} and the two doubling rates.
\omcode{code/methods/29-games-auctions-and-information/python/qm_games.py}{116}{139}{Kelly with and without side information (Python).}
\item \textbf{Run} \texttt{multiplicative\_weights(RPS)}, \texttt{best\_response\_curve()}, \texttt{equivalence()}, \texttt{reserve\_gain(2)}, \texttt{reserve\_gain(5)}, \texttt{treasury\_switch()},
\texttt{winners\_curse()} and \texttt{simulate\_races(0.6)} in \texttt{qm\_games.py}; then \texttt{fig\_games.py}.
\end{tutsteps}
\textbf{What to change next.} Give each dealer demand for two units and find the uniform-price bids by best-response iteration; make the bidders risk averse and watch the first-price revenue overtake the second.
\end{tutorial}

\section{Build: the auction simulator}\label{bld:qm:games-auctions-and-information:bidding}

\begin{build}
\bfield{Purpose} A test bench for the firm's bidding in primary auctions, block trades and request-for-quote competitions: formats, equilibrium benchmarks, \omterm{def:qm:games-auctions-and-information:open}{reserve prices}.
\bfield{Interface} \texttt{fp\_bid}, \texttt{fp\_bid\_general}, \texttt{pab\_bid}; \texttt{simulate(fmt, n, n\_auctions, seed, reserve, increment, values)} for \texttt{first}, \texttt{second},
\texttt{english}, \texttt{dutch}; \texttt{simulate\_multiunit(fmt, n, k, n\_auctions, seed)} for \texttt{uniform}, \texttt{payasbid}; \texttt{expected\_revenue\_uniform}; \texttt{myerson\_reserve}; \texttt{common\_value}.
\bfield{Rules} Every simulated format is checked against its equilibrium benchmark before it is used to evaluate a bidding rule; common-value experiments always report the winners' profit, not the bidders'.
\bfield{Acceptance tests} \texttt{code/firm/bidding/tests/}: bid functions against closed forms (the general first-price formula, pay-as-bid with one unit equal to first price); revenue equivalence in four
formats; reserves for uniform and non-uniform values and a positive seller value; multi-unit revenues and spreads; the winner's curse.
\bfield{Stretch} Multi-unit demand and demand reduction; asymmetric bidders by numerical solution of the equilibrium differential equations; affiliated signals.
\end{build}

\begin{omsources}
\item P.~F.~Malvey and C.~M.~Archibald, ``Uniform-price auctions: update of the Treasury experience'', US Treasury, Office of Market Finance, October 1998; TreasuryDirect, ``How auctions work''.
\item J.~F.~Nash, \emph{Proceedings of the National Academy of Sciences} 36, 1950; J.~von Neumann, \emph{Mathematische Annalen} 100, 1928; J.~C.~Harsanyi, \emph{Management Science} 14, 1967--68.
\item W.~Vickrey, ``Counterspeculation, auctions, and competitive sealed tenders'', \emph{Journal of Finance} 16, 1961; R.~B.~Myerson, ``Optimal auction design'', \emph{Mathematics of Operations Research} 6,
1981; J.~G.~Riley and W.~F.~Samuelson, ``Optimal auctions'', \emph{American Economic Review} 71, 1981; P.~R.~Milgrom and R.~J.~Weber, \emph{Econometrica} 50, 1982.
\item C.~E.~Shannon, \emph{Bell System Technical Journal} 27, 1948; J.~L.~Kelly, ``A new interpretation of information rate'', \emph{Bell System Technical Journal} 35, 1956; T.~M.~Cover and J.~A.~Thomas,
\emph{Elements of Information Theory}, 2nd ed., Wiley, 2006.
\end{omsources}

\section{Exercises}

\begin{exercise}[$\star$]\label{exo:qm:games-auctions-and-information:1}
Ten bidders with uniform private values bid in a \omterm{def:qm:games-auctions-and-information:auctions}{first-price auction}. What does a bidder with value 0.7 bid, and what is the expected revenue?
\end{exercise}

\begin{exercise}[$\star$]\label{exo:qm:games-auctions-and-information:2}
Find the mixed equilibrium of matching pennies with payoffs $\pm1$, and its value.
\end{exercise}

\begin{exercise}[$\star$]\label{exo:qm:games-auctions-and-information:3}
Compute the \omterm{def:qm:games-auctions-and-information:entropy}{entropy} of a fair die, and the \omterm{def:qm:games-auctions-and-information:entropy}{mutual information} between the die and its parity.
\end{exercise}

\begin{exercise}[$\star\star$]\label{exo:qm:games-auctions-and-information:4}
Derive the first-price equilibrium bid for values with distribution function $F(v) = v^2$ on $[0, 1]$, and check it against the general formula.
\end{exercise}

\begin{exercise}[$\star\star$]\label{exo:qm:games-auctions-and-information:5}
Find the optimal reserve for uniform values on $[0, 1]$ when the seller values the object at 0.2.
\end{exercise}

\begin{exercise}[$\star\star$]\label{exo:qm:games-auctions-and-information:6}
Show that in a \omterm{def:qm:games-auctions-and-information:pab}{pay-as-bid auction} of $k$ units to unit-demand bidders, $k = 1$ gives back the first-price equilibrium.
\end{exercise}

\begin{exercise}[$\star\star\star$]\label{exo:qm:games-auctions-and-information:7}
\emph{Coding.} Simulate the four single-unit formats with five bidders, a million auctions each, with and without the optimal reserve, and compare with the closed form.
\end{exercise}

\begin{exercise}[$\star\star\star$]\label{exo:qm:games-auctions-and-information:8}
\emph{Find the flaw.} ``We won 30 of the last 30 block trades we bid for, each at our model's fair value; the model is excellent.''
\end{exercise}

\section{Problem: The Treasury's Switch}

\begin{problem}[{Weekend problem --- pay your bid or pay the clearing price}]\label{pb:qm:games-auctions-and-information:1}
Five dealers bid for securities; the Treasury wonders whether a single clearing price would raise more than pay-as-bid.

\medskip\noindent\textbf{Part I --- Equilibrium.}
\begin{enumerate}
\item What does each dealer bid in a single-unit \omterm{def:qm:games-auctions-and-information:auctions}{first-price auction} with uniform values, and why?
\item Why is truthful bidding dominant in the \omterm{def:qm:games-auctions-and-information:auctions}{second-price auction}?
\item What does the best-response curve of a dealer with value 0.8 show?
\item What does revenue equivalence predict for the four single-unit formats, and what does simulation give?
\item Why do the English and Dutch clocks with ticks of 0.01 miss by half a tick in opposite directions?
\end{enumerate}

\medskip\noindent\textbf{Part II --- Reserve and multi-unit.}
\begin{enumerate}[resume]
\item What is the optimal reserve and what does it gain with two and with five bidders?
\item What do dealers bid in the two-unit \omterm{def:qm:games-auctions-and-information:pab}{pay-as-bid auction}?
\item What revenue do the two multi-unit formats raise?
\item What differs between them, and how does that match the Treasury's study?
\item Why could the model not settle the question the Treasury asked?
\end{enumerate}

\medskip\noindent\textbf{Part III --- Information.}
\begin{enumerate}[resume]
\item How large is the winner's curse in the common-value experiment, and how is it removed?
\item What is the \omterm{def:qm:games-auctions-and-information:entropy}{entropy} of the race, and the doubling rate without a tip?
\item What does a tip right 60\% of the time carry, and what does simulation give?
\item Why do the odds not matter for the gain?
\item What is a tip right 25\% of the time worth, and why?
\end{enumerate}

\medskip\noindent\textbf{Part IV --- Judgement.}
\begin{enumerate}[resume]
\item What would make the uniform-price format raise more in reality?
\item When should the firm shade its bids in a block-trade competition?
\item How much would you pay for a tip, in doubling rate?
\item State the \emph{named result}: the expected revenue of the two formats in the symmetric private-value model, the equilibrium shading of five dealers, and the gain from the optimal reserve.
\item In one sentence: what does revenue equivalence say about auction design?
\end{enumerate}
\end{problem}

\section{Interview questions}

\begin{interviewq}[$\star$ \iqroles{trader, researcher}]\label{iq:qm:games-auctions-and-information:1}
Two bidders with values uniform on $[0, 1]$ bid in a \omterm{def:qm:games-auctions-and-information:auctions}{first-price auction}. What should you bid with value $v$?
\end{interviewq}

\begin{interviewq}[$\star\star$ \iqroles{trader}]\label{iq:qm:games-auctions-and-information:2}
What is the winner's curse and how do you bid to avoid it?
\end{interviewq}

\begin{interviewq}[$\star\star$ \iqroles{researcher}]\label{iq:qm:games-auctions-and-information:3}
State the revenue equivalence theorem and give two ways it fails in practice.
\end{interviewq}

\begin{interviewq}[$\star\star$ \iqroles{researcher, trader}]\label{iq:qm:games-auctions-and-information:4}
How much is a signal worth to a Kelly bettor?
\end{interviewq}

\begin{interviewq}[$\star\star$ \iqroles{researcher}]\label{iq:qm:games-auctions-and-information:5}
Find the value and optimal strategies of the \omterm{def:qm:games-auctions-and-information:game}{zero-sum game} with payoff matrix $\begin{pmatrix}3 & -1\\ -2 & 1\end{pmatrix}$ for the row player.
\end{interviewq}

\begin{interviewq}[$\star\star\star$ \iqroles{trader, researcher}]\label{iq:qm:games-auctions-and-information:6}
You bid for a portfolio in a request-for-quote against four dealers. How does the number of competitors change your bid?
\end{interviewq}
